%% OCCASION = MACS seminar, Montpellier 2026

%% Requires compilation with XeLaTeX or LuaLaTeX
\documentclass[compress,10pt,xcolor={table,dvipsnames},t,aspectratio=169]{beamer}
\usetheme{diapo}

\usepackage{amsmath}
\usepackage{amssymb}
\usepackage{xcolor}
\usepackage{tikz}
\usepackage{etoolbox}

\setlength{\parindent}{0pt}

% bibliographie
\usepackage[style=authoryear,sorting=nyt,doi=false,url=false,maxbibnames=99,date=year]{biblatex}
\addbibresource{biblio.bib}
\setbeamertemplate{bibliography item}{}
\defbibenvironment{bibliography}
  {\list{}{\leftmargin\bibhang\itemindent-\bibhang\itemsep\bibitemsep\parsep\bibparsep}}
  {\endlist}
  {\item[]}
\renewcommand*{\bibfont}{\footnotesize}
\renewrobustcmd*{\bibopenparen}{[}
\renewrobustcmd*{\bibcloseparen}{]}
\renewcommand*{\nameyeardelim}{\addcomma\space}
\newcommand*{\myfootcite}[1]{\footnote[frame,1]{\parencite{#1}}}

\hypersetup{
	colorlinks=true,
	urlcolor=blue,
	citecolor=other,
	linkcolor=title,
}

\title[Hybrid FE and NN methods for parametric PDEs]{Hybrid finite element and neural network methods for parametric PDEs}
\subtitle{\small MACS seminar \\
(``Modélisation, Analyse et Calcul Scientifique'')}

\author[LECOURTIER Frédérique]{%
	Michel Duprez\inst{1},
	Emmanuel Franck\inst{2}, \\
	\textbf{Frédérique Lecourtier}\inst{1} and
	Vanessa Lleras\inst{3}
}

\institute{%
	\inst{1} Project-Team MIMESIS, Inria, Strasbourg \\
	\inst{2} Project-Team MACARON, Inria, Strasbourg \\
	\inst{3} IMAG, University of Montpellier, Montpellier
}

\date{November 10, 2026}

\begin{document}
	\pgfkeys{/pgf/number format/.cd,1000 sep={\,}} % spacing between thousands in numbers

	% page de titre (palette duo : triangles bleu et vert)
	\SetPalette{duo}
	{\setbeamertemplate{footline}{}
		\BackgroundTitle
		\begin{frame}[plain]
			\maketitle
		\end{frame}
	}
	\addtocounter{framenumber}{-1}

	\Background

	%% ---------- Partie commune ----------
	\Chapter*{duo}
	\Section*{Introduction}

\begin{frame}{Context and motivation}
	\todo
\end{frame}

\begin{frame}{Outline}
	\todo
	% \begin{itemize}
	% 	\item[\contribdot{cblue}] Contribution 1 : TODO
	% 	\item[\contribdot{cgreen}] Contribution 2 : TODO
	% \end{itemize}
\end{frame}

\begin{frame}{Elliptic parametric PDEs}
	\todo
\end{frame}
	\section[Classical methods]{Classical numerical methods for PDEs}

\begin{frame}{Discontinuous Galerkin method}
	\todo[Espace discontinu, flux numériques, pénalisation\\Comparaison CG / DG]
\end{frame}

\begin{frame}{Continuous Galerkin method}
	\todo[Problème modèle, formulation faible\\Espace $\mathbb{P}_k$ de Lagrange continu]
\end{frame}

\begin{frame}{Physics-informed neural networks}
	\todo[Principe, fonction de coût\\Version paramétrique]
\end{frame}


	%% ---------- Contribution 1 (bleu) ----------
	\Chapter{blue}{1st hybridization}{Enriched finite element method using PINNs}
	\section{Presentation}

\begin{frame}{Pipeline of the hybrid method}
	\todo
\end{frame}

\begin{frame}{Additive approach}
	\todo
\end{frame}

\begin{frame}{Convergence analysis}
	\todo
\end{frame}

\section{Numerical results}

\subsection{Test case 1}

\begin{frame}{Test case 1}
	\todo
\end{frame}

\subsection{Test case 2}

\begin{frame}{Test case 2}
	\todo
\end{frame}

\section{Conclusion, limitations and perspectives}

\begin{frame}{Conclusion, limitations and perspectives}
	\todo
\end{frame}

	%% ---------- Contribution 2 (vert) ----------
	\Chapter{green}{2nd hybridization}{Learnable Galerkin methods}
	\section{Introduction}

\begin{frame}{Autodiff}
	\todo
\end{frame}

\section{Formalization}

\begin{frame}{Optimization problem}
	\todo
\end{frame}

\section[PINN loss]{Non-parametric enrichment (PINN loss)}

\begin{frame}{Algorithm}
	\todo
\end{frame}

\section[Neural operator]{Parametric enrichment (neural operator)}

\begin{frame}{Algorithm}
	\todo
\end{frame}


	%% ---------- Conclusion ----------
	\Chapter*{duo}
	\section{Conclusion}

\begin{frame}{Conclusion}
	\begin{contribbox}{cblue}{Contribution 1 : TODO}
		TODO
	\end{contribbox}

	\vspace{0.3cm}

	\begin{contribbox}{cgreen}{Contribution 2 : TODO}
		TODO
	\end{contribbox}

	\vspace{0.3cm}

	\textbf{Perspectives :}
	\begin{itemize}
		\item TODO
	\end{itemize}
\end{frame}


	%% ---------- Annexes (une par hybridation) ----------
	\Appendix{blue}{1st hybridization}{Enriched finite element method using PINNs}
	\section{Additional details}
\begin{frame}{Additional details}
	\todo[Détails de la méthode (hors présentation principale)]
\end{frame}

\section[Extra results]{Additional numerical results}
\begin{frame}{Additional numerical results}
	\todo[Résultats numériques supplémentaires]
\end{frame}


	\Appendix{green}{2nd hybridization}{Learnable Galerkin methods}
	\section{Additional details}
\begin{frame}{Additional details}
	\todo[Détails de la méthode (hors présentation principale)]
\end{frame}

\section[Extra results]{Additional numerical results}
\begin{frame}{Additional numerical results}
	\todo[Résultats numériques supplémentaires]
\end{frame}


\end{document}
